\documentclass[11pt,a4paper,onecolumn]{exam}
\usepackage{amsmath}
\usepackage{amssymb}
\usepackage{graphicx}
\usepackage{float}
\usepackage{booktabs}
\usepackage{geometry}
\geometry{top=1.4cm, bottom=1.4cm, left=1.8cm, right=1.8cm}
\usepackage{tikz}
\usepackage[skins]{tcolorbox}

\newcommand{\questionheader}[1]{%
  \begin{tcolorbox}[
    enhanced,
    colback=black,
    coltext=white,
    boxrule=0pt,
    fontupper=\bfseries,
    arc=3mm,
    top=1.5pt, bottom=1.5pt
  ]
  #1
  \end{tcolorbox}%
}

\usepackage{caption}
\usepackage{subcaption}

\newenvironment{blanksolution}
  {%
    \renewcommand{\solutiontitle}{\noindent}%
    \begin{solution}%
  }%
  {\end{solution}}

\pagestyle{headandfoot}
\runningheader{Assignment 7}{Signal Processing in Practice}{Dwaipayan Haldar}

\begin{document}

\begingroup
    \centering
    \LARGE E9 222 Signal Processing in Practice\\
    \LARGE Assignment 7\\[0.5em]
    \large \today\\[0.5em]
    \large Dwaipayan Haldar\par
\endgroup
\noindent\rule{\textwidth}{0.5pt}
\printanswers
\renewcommand{\solutiontitle}{\noindent\textbf{Ans:}\enspace}
\setlength{\parskip}{1pt}

% ─────────────────────────────────────────────────────────────────
\questionheader{Problem 1 (5 Marks): Train ResNet18 on 20 Base Classes}

\noindent ResNet18 (no pretrained weights) was trained for 40 epochs on the first 20 classes of the Oxford-IIIT-Pet dataset using Adam with weight decay and a step LR scheduler.

\vspace{2pt}
\begin{center}
\begin{tabular}{lc}
\toprule
\textbf{Metric} & \textbf{Value} \\
\midrule
Test Accuracy on 20 base classes & 42.20\% \\
Feature dimension $f$ & 512 \\
Average feature norm $\frac{1}{N}\sum_{i=1}^{N}\|F_0(x_i)\|$ & 24.75 \\
\bottomrule
\end{tabular}
\end{center}

% ─────────────────────────────────────────────────────────────────
\questionheader{Problem 2 (10 Marks): Incremental Learning — Expanded Classifier $G_1 \in \mathbb{R}^{37 \times f}$}

\noindent The FC layer is expanded from $20\!\times\!512$ to $37\!\times\!512$, preserving original weights. \textbf{Setting 1} fine-tunes both $F$ and $G_1$ on only the new 17 classes. \textbf{Setting 2} freezes $F_0$ and trains only $G_1$ on new 17 classes and all 37 classes. 
\vspace{2pt}
\begin{center}
\begin{tabular}{lccc}
\toprule
\textbf{Method} & \textbf{Total Acc.} & \textbf{Base 20 Acc.} & \textbf{New 17 Acc.} \\
\midrule
S1: Fine-tune $F + G_1$ (new 17 only) & 23.52\% & 0.00\% & 51.13\% \\
S2: Frozen $F_0$, train $G_1$ (new 17 only) & 17.55\% & 0.00\% & 38.15\% \\
S2: Frozen $F_0$, train $G_1$ (all 37) & 27.53\% & 32.41\% & 21.80\% \\
\bottomrule
\end{tabular}
\end{center}

\vspace{1pt}
\noindent\textbf{Observation:} Setting 1 exhibits catastrophic forgetting — updating the backbone on only new-class data shifts its representations, collapsing base accuracy to 0\%. Setting 2 with only new data also exhibits the same effect with the base 20 class accuracy going to 0. Setting 2 if done will all the classes preserves base accuracy (32.86\%) by keeping $F_0$ frozen, but new-class accuracy is limited (20.62\%). This shows that replaying the old data is the only way to prevent catastrophic forgetting. Next part of the question shows even without replaying all the old data, information and accuracy can be preserved.

% ─────────────────────────────────────────────────────────────────
\questionheader{Problem 3 (15 Marks): Mitigating Catastrophic Forgetting via Feature Distribution Replay}

\noindent For each of the 20 base classes, the mean $\mu_c$ and full covariance $\Sigma_c$ of $F_0$ embeddings are computed. Five synthetic features per class are sampled as $\tilde{f} \sim \mathcal{N}(\mu_c,\,\Sigma_c + \epsilon I)$ ($\epsilon=10^{-4}$, 100 synthetic old-class samples total). $G_1$ is trained on these plus 1,695 real new-class feature vectors, with $F_0$ frozen.

\vspace{2pt}
\begin{center}
\begin{tabular}{lccc}
\toprule
\textbf{Method} & \textbf{Total Acc.} & \textbf{Base 20 Acc.} & \textbf{New 17 Acc.} \\
\midrule
P2-S1: Fine-tune $F+G_1$ (new only)        & 23.52\% & 0.00\%  & 51.13\% \\
P2-S2: Frozen $F_0$, train $G_1$ (all 37)  & 27.23\% & 32.86\% & 20.62\% \\
P3:\phantom{-S2} Frozen $F_0$ + replay     & 21.31\% & 5.50\%  & 39.87\% \\
\bottomrule
\end{tabular}
\end{center}

\vspace{1pt}
\noindent\textbf{Observation:} Using a full covariance Gaussian (capturing inter-feature correlations) improves base accuracy over a diagonal approximation (5.50\% vs 1.11\%). Replay reduces forgetting vs S1 (base: 5.50\% vs 0.00\%) and yields higher new-class accuracy than S2 (39.87\% vs 20.62\%). However, base accuracy still lags S2 (32.86\%) due to class imbalance: 100 synthetic old-class samples against 1,695 real new-class samples biases the classifier toward new classes.

\noindent\textbf{Why feature statistics suffice:} A trained backbone maps images to compact, near-Gaussian clusters in feature space. Storing only $\mu_c$ and $\Sigma_c$ per class — without any raw images — captures the class geometry well enough for a linear classifier to remain discriminative. This makes feature distribution replay both memory-efficient and privacy-preserving, while meaningfully mitigating forgetting.

\end{document}
